% =====================================================================
% Packages
% =====================================================================
\usepackage[utf8]{inputenc}
\usepackage[T1]{fontenc}
\usepackage{lmodern}
\usepackage{microtype}

\usepackage{amsmath,amssymb,amsthm,mathtools}
\usepackage{booktabs}
\usepackage{array}
\usepackage{tabularx}
\usepackage{enumitem}

\usepackage[colorlinks=true,linkcolor=blue!50!black,citecolor=blue!50!black,urlcolor=blue!50!black]{hyperref}
\usepackage[capitalise,noabbrev]{cleveref}
\usepackage{url}

\usepackage{xcolor}
% Skeleton placeholders. \stub{...} is a robust paragraph-level marker
% that handles itemize/enumerate/math display inside; do NOT use the
% todonotes package — amsart's narrow marginpar geometry breaks its tikz
% rendering.
% For final submission, redefine \stub to swallow its argument; see
% Makefile target `final`.
\providecommand{\stub}[1]{%
  \par\medskip\noindent
  \fcolorbox{orange!60!black}{orange!8}{%
    \begin{minipage}{\dimexpr\linewidth-2\fboxsep-2\fboxrule\relax}
    \small\itshape\textbf{[TODO]}\hspace{0.5em}#1
    \end{minipage}%
  }\par\medskip
}
% Alias \todo to \stub for compatibility with any leftover macros.
\providecommand{\todo}[1]{\stub{#1}}

% =====================================================================
% Theorem environments (AMS style)
% =====================================================================
\theoremstyle{plain}
\newtheorem{theorem}{Theorem}[section]
\newtheorem{lemma}[theorem]{Lemma}
\newtheorem{proposition}[theorem]{Proposition}
\newtheorem{corollary}[theorem]{Corollary}

\theoremstyle{definition}
\newtheorem{definition}[theorem]{Definition}
\newtheorem{example}[theorem]{Example}
\newtheorem{problem}{Problem}[section]

\theoremstyle{remark}
\newtheorem{remark}[theorem]{Remark}
\newtheorem{conjecture}[theorem]{Conjecture}

% =====================================================================
% Macros — graphs and basic objects
% =====================================================================
\newcommand{\Adj}{A}                          % adjacency matrix
\newcommand{\graph}{G}
\newcommand{\Hgraph}{H}
\newcommand{\splus}{s^{+}}                    % positive square energy
\newcommand{\sminus}{s^{-}}                   % negative square energy
\newcommand{\splusG}{\splus(\graph)}
\newcommand{\sminusG}{\sminus(\graph)}
\newcommand{\splusH}{\splus(\Hgraph)}
\newcommand{\sminusH}{\sminus(\Hgraph)}

% Ear-gain quantities
\newcommand{\dplus}{\delta^{+}}
\newcommand{\dminus}{\delta^{-}}

% Two-tree families
\newcommand{\Bk}{B_{k}}
\newcommand{\Ln}{L_{n}}
\newcommand{\BTkt}[2]{\mathrm{BT}(#1,#2)}
\newcommand{\Fn}{F_{n}}
\newcommand{\Spider}[3]{S(#1,#2,#3)}

% =====================================================================
% Macros — spectral quantities at an ear
% =====================================================================
\newcommand{\earw}{w}                         % w = e_a + e_b
\newcommand{\earc}{c}                         % c_i = w^T u_i
\newcommand{\earmu}{\mu}                      % eigenvalue of A(H)
\newcommand{\Wm}{W^{-}}
\newcommand{\Wp}{W^{+}}
\newcommand{\Wz}{W^{0}}
\newcommand{\Mkm}{M_{k}^{-}}
\newcommand{\Mkp}{M_{k}^{+}}
\newcommand{\Mone}{M_{1}}
\newcommand{\Mtwo}{M_{2}}
% Optional argument lets these compose with extra subscripts (e.g.
% asymptotic forms): \Monem -> M_1^-, \Monem[1,\infty] -> M_{1,\infty}^-.
\newcommand{\Monem}[1][1]{M_{#1}^{-}}
\newcommand{\Mtwom}[1][2]{M_{#1}^{-}}
\newcommand{\Monep}[1][1]{M_{#1}^{+}}
\newcommand{\Mtwop}[1][2]{M_{#1}^{+}}

% Joint invariant and threshold
\newcommand{\Iv}{I(v)}
\newcommand{\Ivs}{I(v^{*})}
\newcommand{\Iinf}[1]{I_{\infty}(#1)}
\newcommand{\threshold}{T}

% Selector quantities
\newcommand{\degsum}{\sigma}                  % deg_H(a) + deg_H(b)
\newcommand{\cliqtree}{T}                     % clique tree of G
\newcommand{\maxdeg}{v^{*}}                   % max-degsum ear

% =====================================================================
% Conventions / phrases
% =====================================================================
\newcommand{\Lprime}{(\mathrm{L}')}
\DeclareMathOperator{\spec}{spec}
\DeclareMathOperator{\tr}{tr}
\DeclareMathOperator{\rank}{rank}
\DeclareMathOperator{\diag}{diag}
\DeclareMathOperator*{\argmax}{arg\,max}
\DeclareMathOperator*{\argmin}{arg\,min}
\DeclareMathOperator{\sgn}{sgn}
\DeclareMathOperator{\diam}{diam}

% Reference shortcuts
\newcommand{\arXiv}[1]{\href{https://arxiv.org/abs/#1}{arXiv:#1}}
